\subsection{Міжнародний бенчмаркінг: українська дизайн-освіта у глобальному контексті}\label{subsec:2.3}

Для контекстуалізації результатів загальнонаціонального аналізу (п.~\ref{subsec:2.1}) та кейс-стаді КДПУ (п.~\ref{subsec:2.2}) проведено порівняльний аналіз інтеграції ШІ/ІКК-змісту у провідних міжнародних програмах дизайну. Вибірка охоплює 8 програм із 6 країн, що представляють різні моделі дизайн-освіти.

\subsubsection{Методика бенчмаркінгу}\label{sssec:2.3.1}

Для порівняння використано уніфіковану рамку, що відображає 4 компоненти ІКК:
\begin{enumerate}
    \item \textbf{ІА}~--- інформаційно-аналітичний: наявність дисциплін із аналізу даних, дослідницьких методів для дизайнерів, критичної оцінки ШІ-контенту;
    \item \textbf{К}~--- комунікативний: цифрова презентація, портфоліо, командна робота з використанням цифрових платформ;
    \item \textbf{Т}~--- технологічний: ШІ-інструменти, генеративний дизайн, обчислювальний дизайн, промпт-інженерія;
    \item \textbf{Р}~--- рефлексивний: етика ШІ в дизайні, критичний дизайн, дизайн-мислення в контексті ШІ.
\end{enumerate}

Дані отримано з відкритих джерел: офіційних веб-сайтів університетів, каталогів курсів та публікацій про оновлення програм.

\subsubsection{Результати порівняння}\label{sssec:2.3.2}

\begin{longtable}{@{}p{0.25\textwidth}p{0.18\textwidth}p{0.08\textwidth}p{0.08\textwidth}p{0.08\textwidth}p{0.08\textwidth}p{0.10\textwidth}@{}}
\caption{Покриття компонентів ІКК у міжнародних програмах дизайну}\label{tab:international-benchmark} \\
\hline
\textbf{Університет / Країна} & \textbf{Програма} & \textbf{ІА} & \textbf{К} & \textbf{Т} & \textbf{Р} & \textbf{ШІ-дисципліни} \\
\hline
\endfirsthead
\hline
\textbf{Університет / Країна} & \textbf{Програма} & \textbf{ІА} & \textbf{К} & \textbf{Т} & \textbf{Р} & \textbf{ШІ-дисципліни} \\
\hline
\endhead
Royal College of Art, Великобританія & MA Design & +++ & +++ & +++ & +++ & 3--4 \\
Parsons School, США & BFA Design \& Technology & ++ & +++ & +++ & ++ & 2--3 \\
Aalto University, Фінляндія & BA Design & ++ & ++ & +++ & +++ & 2--3 \\
TU Delft, Нідерланди & BSc Industrial Design & +++ & ++ & ++ & ++ & 1--2 \\
RMIT, Австралія & BA Design & ++ & +++ & ++ & ++ & 1--2 \\
Politecnico di Milano, Італія & BSc Design & ++ & ++ & ++ & + & 1 \\
Seoul National, Південна Корея & BFA Design & ++ & ++ & +++ & + & 2--3 \\
\addlinespace
\textbf{КДПУ (середнє)} & \textbf{3 програми} & \textbf{+} & \textbf{+} & \textbf{+} & \textbf{+/--} & \textbf{0} \\
\addlinespace
\textbf{Україна (середнє)} & \textbf{122 прогр.} & \textbf{+} & \textbf{+} & \textbf{+} & \textbf{+/--} & \textbf{0} \\
\hline
\end{longtable}

\textit{Примітка:} +++ = повна інтеграція (окрема дисципліна + наскрізна інтеграція); ++ = часткова (окремий модуль або наскрізні елементи); + = мінімальна (окремі згадки); +/-- = відсутня або формальна.

\subsubsection{Ключові спостереження}\label{sssec:2.3.3}

\textbf{1. Спеціалізовані ШІ-дисципліни.} Усі 7 міжнародних програм мають від 1 до 4 дисциплін, безпосередньо присвячених ШІ в дизайні. Приклади: <<AI for Design>> (RCA), <<Generative Systems>> (Parsons), <<Computational Design>> (Aalto), <<Design \& AI Ethics>> (TU Delft). Жодна з 122 українських програм не має аналогічних дисциплін.

\textbf{2. Промпт-інженерія як формальна дисципліна.} Провідні школи (Georgia Tech, Columbia, RISD) включають промпт-інженерію у навчальні плани як самостійний модуль або компонент існуючих курсів~\cite{Musiienko2026ScopingReview}. Цей аспект технологічного компоненту ІКК повністю відсутній в українських програмах.

\textbf{3. Рефлексивний компонент~--- найбільший розрив.} Міжнародні програми активно інтегрують курси з етики ШІ, критичного дизайну та соціальної відповідальності дизайнера. В українських програмах рефлексивний компонент (самооцінка, етика, критичний аналіз ШІ) практично не представлений.

\textbf{4. Наскрізна інтеграція vs. ізольовані дисципліни.} Найуспішніші міжнародні програми (RCA, Aalto) інтегрують ШІ-тематику \textit{наскрізно} через усі дисципліни (від першого до останнього курсу), а не лише через окремі курси. Це є принциповим підходом для моделі трансформації (розділ~\ref{sec:3}).

\textbf{5. Географічна кореляція.} Аналіз оглядового дослідження~\cite{Musiienko2026ScopingReview} виявив, що країни-лідери досліджень генеративного ШІ в дизайн-освіті (Китай~--- 31,4\%, США~--- 13,5\%, Великобританія~--- 4,5\%) також лідирують у практичній інтеграції ШІ у навчальні програми. Україна не представлена у жодному з 156 досліджень, що підтверджує відставання.

\subsubsection{Найкращі практики для адаптації}\label{sssec:2.3.4}

На основі міжнародного бенчмаркінгу визначено 5 найкращих практик, придатних для адаптації в українських програмах:

\begin{enumerate}
    \item \textbf{Прогресивна інтеграція ШІ}: від базових модулів ШІ-грамотності (1 курс) через використання ШІ-інструментів у фахових дисциплінах (2--3 курси) до критичного та етичного осмислення (4 курс)~--- модель Aalto University;
    \item \textbf{Паралельне збереження ручних навичок}: жодна міжнародна програма не замінює рисунок, живопис чи макетування ШІ-інструментами. ШІ \textit{доповнює}, а не заміщує традиційні навички~--- підхід RCA;
    \item \textbf{Промпт-інженерія як структурована компетентність}: від <<основ промптингу>> через <<ланцюговий промптинг>> до <<критичної оцінки ШІ-виходів>>~--- підхід Parsons/Georgia Tech;
    \item \textbf{Етика та критичний аналіз ШІ}: окремий курс або модуль з ШІ-етики (bias, авторське право, екологічний вплив)~--- підхід TU Delft;
    \item \textbf{Портфоліо-орієнтована оцінка}: міжнародні програми оцінюють ІКК через цифрове портфоліо, що демонструє як ШІ-проєкти, так і традиційні роботи~--- підхід RMIT.
\end{enumerate}

Ці практики враховані при розробці моделі трансформації навчальних програм КДПУ (розділ~\ref{sec:3}).
